\section{Experiments}\label{sec:experiments}

Ten findings, eight core calibrations plus two reviewer-motivated probes, and three follow-ups that close questions the ten left open: an existence witness for rank-eligibility (\cref{sec:witness}), an asymmetric positive control for the manipulation probe (\cref{sec:f5-positive-control}) and sealed evaluation seeds (\cref{sec:sealed-seeds}). Each subsection states the entry point, seeds and quantity measured, then reports numbers from committed JSON and figures. The exact commands are listed in \cref{app:commands}. The original evidence records SharpeBench 0.5.0 under \texttt{sharpearena} 0.8.0; the concave-impact and predictability extensions were produced on the v0.11.0 source tree, and the three follow-ups on the post-0.11.1 tree, each with its configuration serialized beside its results.

Three scoping statements govern the whole section. First, the findings are instrument calibrations on a generator whose canonical tape the realism gate of \cref{sec:f4} declines to certify (1 of 24 runs passes the full conjunction); certification is reported beside the results, not required by them. Second, the byte-identity guarantee covers engine outputs behind F1 through F3 and the canonical linear path; Python evidence reductions are seed-deterministic on one platform but outside the golden-hash path. Third, the core evidence uses the explicitly stated train-band seeds except F3, which exercises the split; the predictability probe uses a separate 50{,}000-start band and reports why a bounded public band is unsafe. Dispersion is reported throughout: F1 and F3 use 2{,}000 seed bootstrap resamples, while F2, F5 and F6 use t-based 95\% intervals over committed vectors.

\subsection{F1: Pipeline validation under the corrected kernel}\label{sec:f1}

The reference-policy field (flat, equal-weight long, momentum, minimum-variance, maximum-Sharpe, fractional-Kelly volatility targeting) is rolled over sixteen scenario seeds in each of the Calm, Hard and Extreme tiers via \texttt{run\_baselines}. Each policy's pooled return series is scored by the SharpeBench \texttt{score\_run} kernel with the declared trial count equal to the field size, and each row carries a seed-paired bootstrap 95\% confidence interval on the deflated Sharpe (2{,}000 resamples).

This experiment validates the producer-scorer pipeline end to end; the scorer-side discovery it reproduces belongs to the companion paper. Under the pre-0.5.0 kernel this board was all zeros: every baseline scored a deflated Sharpe of 0.0000 on every tier. That was not strictness. The kernel applied its annualized 0.5 deflation prior per period without conversion, which set the bar near an annualized Sharpe of 18 on daily bars, a level nothing clears. Finding and fixing that unit bug is the companion SharpeBench paper's Finding 1 \citep{toca2026sharpebench} (the derivation: an annualized prior applied per period, corrected by converting once at 252 periods per year); what this paper claims is only that the producer-scorer separation surfaced it, because trajectories produced here could be re-scored under both kernels and the discrepancy attributed to the scorer alone.

\begin{table}[t]
\centering
\caption{Baseline leaderboard under the corrected kernel: deflated Sharpe (DSR) and pass\textsuperscript{$k$} rate per tier, 16 seeds each. No policy is rank-eligible on any tier: eligibility requires the per-run gate to pass on every seed (a pass\textsuperscript{$k$} rate of exactly 1.00), and the best rate is 0.75.}
\label{tab:f1}
\small
\setlength{\tabcolsep}{4.5pt}
\begin{tabular}{l rr rr rr}
\toprule
 & \multicolumn{2}{c}{Calm} & \multicolumn{2}{c}{Hard} & \multicolumn{2}{c}{Extreme} \\
\cmidrule(lr){2-3}\cmidrule(lr){4-5}\cmidrule(lr){6-7}
Policy & DSR & pass\textsuperscript{$k$} & DSR & pass\textsuperscript{$k$} & DSR & pass\textsuperscript{$k$} \\
\midrule
kelly\_vol\_target & 1.0000 & 0.75 & 0.0277 & 0.31 & 0.0029 & 0.12 \\
equal\_weight\_long & 1.0000 & 0.62 & 0.1114 & 0.31 & 0.0234 & 0.06 \\
min\_variance & 0.9849 & 0.44 & 0.0024 & 0.19 & 0.0243 & 0.06 \\
flat & 0.0007 & 0.00 & 0.0007 & 0.00 & 0.0007 & 0.00 \\
momentum & 0.0000 & 0.00 & 0.0000 & 0.00 & 0.0000 & 0.06 \\
max\_sharpe & 0.0000 & 0.06 & 0.0000 & 0.00 & 0.0000 & 0.00 \\
\bottomrule
\end{tabular}
\end{table}

\Cref{tab:f1} and \cref{fig:f1} show the corrected board. On Calm, drift is free deflated Sharpe: \texttt{kelly\_vol\_target} and \texttt{equal\_weight\_long} saturate the DSR at 1.0000 (bootstrap CIs $[0.9798, 1.0000]$ and $[0.8343, 1.0000]$) and \texttt{min\_variance} reaches 0.9849, though with a near-vacuous interval of $[0.0264, 1.0000]$. Yet nothing on the board is rank-eligible, because no policy passes the per-run gate on every seed; the best pass\textsuperscript{$k$} rate is 0.75. On Hard the best DSR collapses to 0.1114 and on Extreme to 0.0243, and the long baseline's pass\textsuperscript{$k$} rate degrades monotonically with stress (0.62, 0.31, 0.06). \texttt{momentum} and \texttt{max\_sharpe} lose money outright on every tier (mean per-bar returns down to $-1.6 \times 10^{-3}$ on Extreme).

The reading: eligibility is a conjunction, and each leg catches what the other misses. Deflation alone would crown drift on Calm; pass\textsuperscript{$k$} alone would miss the overfit breadth the deflation discounts. The earlier all-zeros board looked strict but was measuring a unit error. The corrected board is the honest zero: the number a real agent has to beat is a positive, deflated Sharpe that survives every evaluation seed, and no reference policy produces one. One caveat the board cannot answer on its own: no reference policy is rank-eligible, so the board alone does not show that the eligibility set is non-empty for this task family. \Cref{sec:witness} answers that by construction, injecting an out-of-band oracle signal of controlled strength and locating the strength at which the gates first open; the region is non-empty on every tier, and the per-run gate that fails every baseline here is the leg that binds there too.

\begin{figure}[t]
\centering
\includegraphics[width=\linewidth]{f1-baselines}
\caption{Deflated Sharpe per baseline per tier with seed-paired bootstrap CIs. Calm saturates the long baselines; Hard and Extreme collapse them toward zero.}
\label{fig:f1}
\end{figure}

\subsection{Rank-eligibility witness: the acceptance set is non-empty}\label{sec:witness}

The baseline board of \cref{sec:f1} contains no rank-eligible policy, leaving open whether the conjunction can be satisfied at all. We answer by construction with an out-of-band oracle, not an entrant. For each seed, the script replays the deterministic, policy-invariant tape one bar ahead and mixes standardized next-bar truth $z_t$ with one fixed seeded noise path $\varepsilon_t$:
\[
\mathrm{signal}_t(s)=s z_t+\sqrt{1-s^2}\,\varepsilon_t,
\qquad 0\le s\le1.
\]
The same $\varepsilon_t$ is reused at every $s$ (common random numbers), so changing signal strength does not silently change the experiment. A sign-following policy trades every bar; a deadband policy trades only when $|\mathrm{signal}|>1$ and otherwise holds. Both are scored through the unchanged SharpeBench conjunction over 16 seeds. The script first checks the full coarse acceptance set. It refines a crossing by bisection only when coarse eligibility is monotone; otherwise it reports the observed set without calling any crossing a threshold. All ten attained series are monotone after the common-random-number correction.

\begin{table}[t]
\centering
\caption{Observed eligibility thresholds in signal-truth correlation $s$. Intervals are bisection brackets after a monotone coarse-grid check. ``None'' means that no sampled $s\le1$ is eligible. At every attained threshold, the last failing gate is the requirement that every seed's marginal PSR reach $0.90$.}
\label{tab:witness}
\small
\begin{tabular}{llccc}
\toprule
Variant & Band & Calm & Hard & Extreme \\
\midrule
\texttt{sign\_follow} & held-out & none & $[0.7250,0.7281]$ & $[0.3000,0.3031]$ \\
\texttt{sign\_follow} & F1 table & none & $[0.7875,0.7906]$ & $[0.4438,0.4469]$ \\
\texttt{deadband\_hold} & held-out & $[0.9000,0.9031]$ & $[0.5031,0.5063]$ & $[0.3813,0.3844]$ \\
\texttt{deadband\_hold} & F1 table & $[0.6563,0.6594]$ & $[0.6188,0.6219]$ & $[0.3688,0.3719]$ \\
\bottomrule
\end{tabular}
\end{table}

The result is an existence proof and a calibration, not evidence about an agent. The deadband construction is eligible on every tier, so the acceptance set is non-empty. The every-seed gate is last to open at all ten observed thresholds; pooled DSR and bootstrap gates open earlier. The sign follower never becomes eligible on Calm, even at $s=1$, because its every-bar turnover costs more than the low-volatility signal is worth. Thus the environment does not impose an activity requirement, but its cost model does require a Calm policy to ration turnover.

\begin{figure}[t]
\centering
\includegraphics[width=0.95\linewidth]{witness}
\caption{Witness sweep on the held-out band. Filled markers are eligible sampled points; dashed lines mark bisection thresholds only after monotonicity is verified. The right panels show how many of the 16 seed-level PSR tests pass.}
\label{fig:witness}
\end{figure}


\subsection{F2: Regret against the closed-form reference policy}\label{sec:f2}

On the Avellaneda-Stoikov market-making task, \texttt{mm\_regret} runs each candidate quoting policy and the closed-form reference policy (\cref{sec:tasks}) on the same sixteen seeded episodes ($\sigma = 2.0$, $\gamma = 0.1$, $\kappa = 1.5$, 200 steps) and reports the mean reward gap to the reference. Candidates are the reference itself and fixed-spread quoters over a grid of half-spreads. The evidence commits the per-episode regret vector per grid point; intervals are t-based 95\% CIs over the sixteen paired episodes.

The reference's regret is 0.0 by construction, which confirms the plumbing: the metric's zero point is the reference, not a claim of unbeatability. The fixed-spread curve is U-shaped (\cref{fig:f2}): regret is 84.6 (95\% CI $[68.2, 101.1]$) at a half-spread of 0.05 (quoting too tight, filled constantly, run over by inventory), falls to a minimum of 0.037 at 0.5, and rises again to 36.6 ($[33.0, 40.2]$) at 2.0 and 58.8 ($[55.3, 62.3]$) at 4.0 (quoting too wide, earning nothing). The minimum deserves its interval: the 0.5 point's CI is $[-8.7, 8.8]$ and the 1.0 point's is $[-5.4, 12.4]$, so the best fixed spreads are statistically indistinguishable from the reference on sixteen episodes, while the tight and wide extremes are separated from it by wide margins. The instrument resolves the shape of the curve, not the sign of its minimum.

The point of shipping the reference is that agents on this task are scored on regret, not raw PnL. A market-making leaderboard ranked on PnL rewards whichever policy drew the friendliest flow; regret against a fixed analytical reference on the same episodes removes that luck axis.

\begin{figure}[t]
\centering
\includegraphics[width=0.7\linewidth]{f2-regret}
\caption{Mean regret versus the closed-form Avellaneda-Stoikov reference policy over 16 paired episodes, with t-based 95\% CIs. Fixed-spread quoting is U-shaped in the half-spread; the reference defines zero.}
\label{fig:f2}
\end{figure}

\subsection{F3: Generalization gap and cross-regime transfer}\label{sec:f3}

\texttt{generalization\_gap} scores the default reference policy (equal-weight long) on disjoint train and held-out seed bands (16 seeds each, separated by a 10{,}000-seed gap) within each tier. \texttt{cross\_regime\_transfer} then holds sixteen seeds fixed and varies the regime, producing the full three-by-three transfer matrix whose diagonal is zero by construction, because identical regimes reuse byte-identical environments. One framing note before the numbers: the policy is fixed and parameter-free, so nothing is trained, "source regime" below means the regime the policy is nominally selected on, and for such a policy the matrix reduces to pairwise differences of three per-tier deflated Sharpes. What F3 measures is the instrument, how much score a regime switch moves for a fixed reference, not transfer learning. The evidence commits per-seed return series per cell; CIs are seed-resampled percentile bootstrap ($B = 2000$), seed-paired across regimes for the matrix cells.

\begin{table}[t]
\centering
\caption{Left: within-tier generalization gap (train minus test deflated Sharpe) for the fixed equal-weight reference. Right: cross-regime transfer gap (in-distribution minus zero-shot out-of-distribution DSR), rows source regime, columns scored regime; the diagonal is zero by construction, and the Hard$\leftrightarrow$Extreme cells ($\pm 0.096$) sit below the 0.34 seed-noise calibration and are not interpreted.}
\label{tab:f3}
\small
\begin{tabular}{lrrr}
\toprule
Tier & Train DSR & Test DSR & Gap \\
\midrule
Calm & 1.0000 & 0.9988 & $+0.0012$ \\
Hard & 0.1231 & 0.4618 & $-0.3388$ \\
Extreme & 0.0269 & 0.1280 & $-0.1012$ \\
\bottomrule
\end{tabular}
\hspace{1.5em}
\begin{tabular}{lrrr}
\toprule
 & Calm & Hard & Extreme \\
\midrule
Calm & 0.000 & $+0.877$ & $+0.973$ \\
Hard & $-0.877$ & 0.000 & $+0.096$ \\
Extreme & $-0.973$ & $-0.096$ & 0.000 \\
\bottomrule
\end{tabular}
\end{table}

\Cref{tab:f3} reports both. The within-tier gaps are the expected control: the reference policy has no fitted parameters, so it cannot overfit its train band, and the instrument correctly reads near zero on Calm ($+0.0012$) and noise-signed values on the stressed tiers (the negative Hard and Extreme gaps mean the held-out band happened to be kinder; with no learning, the gap is pure seed noise, and its magnitude, up to 0.34 on Hard, calibrates how much band luck a 16-seed evaluation carries). The per-band bootstrap intervals say the same thing at interval width: Hard's train-band DSR of 0.1231 carries a CI of $[0.000, 0.944]$ and its test band $[0.037, 0.948]$ (the bands are disjoint seed sets, so these are unpaired), overlapping almost entirely, which is what a pure-noise gap looks like.

The transfer matrix (\cref{fig:f3}) is where regime structure appears, and its cells split by the paper's own noise calibration of 0.34. The Calm-involved point estimates are far above it: source regime Calm scored zero-shot on Hard costs $0.877$ of deflated Sharpe, and Calm to Extreme costs $0.973$, essentially the whole score. The seed-paired bootstrap intervals temper the two claims differently: Calm to Extreme excludes zero decisively ($[0.635, 0.999]$), while Calm to Hard's interval is wide ($[0.054, 0.999]$), positive throughout but reaching close to zero, so at 16 seeds only the Calm-to-Extreme collapse is resolved beyond doubt. The Hard-to-Extreme cell ($0.096$, CI $[-0.218, 0.911]$) sits \emph{below} the 0.34 seed-noise calibration and its interval straddles zero, so it is not interpreted: at 16 seeds the instrument cannot distinguish it from band luck. The wide intervals are themselves a calibration result: a bounded, heavily deflated score bootstrapped over 16 seeds is a coarse instrument, and claims narrower than these intervals should not be made from boards this size. The matrix is antisymmetric because the policy is fixed, and its diagonal is zero by construction. Its off-diagonal magnitude is the quantity a within-tier gap is structurally blind to: a policy scored only on calm seeds can carry a saturated DSR into a leaderboard while being worth 0.03 under stress. Both metrics are cheap, and the protocol requires reporting both.

\begin{figure}[t]
\centering
\includegraphics[width=0.62\linewidth]{f3-transfer-matrix}
\caption{Cross-regime transfer matrix for the reference policy: out-of-distribution deflated Sharpe with the seed band held fixed.}
\label{fig:f3}
\end{figure}

\subsection{F4: Stylized-facts realism certification}\label{sec:f4}

For each tier, price panels from eight seeded episodes are graded by \texttt{certify\_realism} against the directional bounds: positive excess kurtosis, positive absolute-return autocorrelation and positive aggregational Gaussianity from the classical catalogue \citep{cont2001empirical}, and a super-Poisson Fano factor (the authors' own diagnostic, \cref{sec:probes}), with gain/loss skew \citep{cont2001empirical} and timescale asymmetry \citep{zumbach2009time} reported as informational.

\begin{table}[t]
\centering
\caption{Stylized-facts certification per tier: mean fact value over 8 seeds, per-fact pass counts, and the all-facts pass rate. Calm is near-Gaussian by design and fails the fat-tail facts; the stressed tiers pass them.}
\label{tab:f4}
\small
\setlength{\tabcolsep}{5pt}
\begin{tabular}{l rr rr rr r}
\toprule
 & \multicolumn{2}{c}{Excess kurtosis} & \multicolumn{2}{c}{$|r|$ autocorr} & \multicolumn{2}{c}{Fano factor} & All facts \\
\cmidrule(lr){2-3}\cmidrule(lr){4-5}\cmidrule(lr){6-7}
Tier & mean & pass & mean & pass & mean & pass & pass rate \\
\midrule
Calm & $-0.97$ & 0/8 & $+0.0002$ & 5/8 & 1.19 & 7/8 & 0.000 \\
Hard & $+16.30$ & 8/8 & $-0.0107$ & 1/8 & 0.85 & 0/8 & 0.000 \\
Extreme & $+11.38$ & 8/8 & $-0.0118$ & 1/8 & 0.79 & 1/8 & 0.125 \\
\bottomrule
\end{tabular}
\end{table}

\Cref{tab:f4} and \cref{fig:f4} report the certification honestly, and it is a mixed verdict. The certification bounds are universal: they encode what real-market tape looks like, not what each tier intends, so a regime designed to be near-Gaussian is expected to fail them, and does. Calm's returns are platykurtic (mean excess kurtosis $-0.97$, aggregational Gaussianity $-0.58$); Calm is therefore, plainly, not realistic tape under this battery, by design, and the gate says so rather than grading it on a curve. Hard and Extreme pass excess kurtosis and aggregational Gaussianity on all eight seeds (means $+16.3$ and $+11.4$ excess kurtosis). But the full conjunction almost never passes: 0/8 on Calm, 0/8 on Hard, 1/8 on Extreme. The binding facts on the stressed tiers are volatility clustering (absolute-return autocorrelation is slightly negative at this 121-bar episode length) and the Fano factor, which drops sub-Poisson exactly when jumps concentrate variance in few bars. This is a finding against the generator, stated as such: the vol-and-jump construction buys fat tails without buying persistent volatility, and a future generator revision must add clustering (for example a stochastic-volatility driver) before the certification gate passes at tier level.

That driver now exists as an opt-in \texttt{vol\_clustering} knob on the scenario spec: when positive, a deterministic EMA persistence state driven by realized absolute returns modulates each bar's return deviation, using the same mul/add/div/max arithmetic as the tier transforms so cross-runtime byte-identity is preserved, and a committed golden hash pins the clustered output alongside the unclustered tiers. At \texttt{vol\_clustering} $= 0$ (the default and the canonical leaderboard configuration) the generated panels are byte-identical to the tapes certified above, so the finding stands as reported, and every other experiment in this section runs on this uncertified default tape. Re-certifying the same tiers and seeds at strength $0.5$ flips the volatility-clustering fact: the absolute-return autocorrelation pass count rises from 5/8 to 8/8 on Calm, 1/8 to 8/8 on Hard, and 1/8 to 5/8 on Extreme (Hard's mean $|r|$ autocorrelation moves from $-0.011$ to $+0.029$), lifting Hard's all-facts pass rate from 0.000 to 0.250. Two honesty notes on that remediation. The strength $0.5$ was tuned, and re-certified, on the same eight seeds used to diagnose the failure, so the improvement is in-sample on the diagnosis panel; certification on a disjoint seed set at several strengths is the required follow-up before the driver can claim the fact generally. And the remaining binding fact on the stressed tiers is the sub-Poisson Fano factor, which the clustering driver improves (Hard 0.85 to 0.90) but does not yet fix.

\begin{figure}[t]
\centering
\includegraphics[width=\linewidth]{f4-realism}
\caption{Per-seed stylized-fact values by tier against the certification bounds.}
\label{fig:f4}
\end{figure}

\subsection{F5: The manipulation boundary and size response}\label{sec:f5}

A push-and-unwind manipulator is scored against its zero-impact paired reference (identical seed and schedule, $\lambda = \eta = 0$) over eight seeds, so \texttt{impact\_pnl} attributes profit specifically to having moved the price. \texttt{impact\_boundary\_sweep} sweeps permanent impact (Kyle $\lambda$), temporary impact ($\eta$) and the momentum-follower gain over six grid points per axis; \texttt{size\_response} sweeps the push weight over five points. The evidence commits per-seed impact-PnL vectors at every grid point; intervals below are t-based 95\% CIs over the eight seeds.

One theoretical fact frames the whole finding. The market composes linear permanent impact with temporary impact, and under linear, time-independent permanent impact, round-trip price manipulation is unprofitable essentially by theorem: linearity is exactly the condition that rules out price-manipulation quasi-arbitrage \citep{huberman2004price}, and the no-dynamic-arbitrage analysis reaches the same conclusion for related impact-and-decay families \citep{gatheral2010no}. The clean sweep below is therefore a consistency check, not a certification that manipulation cannot pay in synthetic markets. \Cref{sec:f5-concave} runs the concave-impact ablation on the same symmetric schedule and reports a null; \cref{sec:f5-positive-control} then searches the asymmetric schedules that null did not cover and finds the profit theory predicts, under concave impact only.

\begin{table}[t]
\centering
\caption{Manipulation probe, impact-attributable PnL as fraction of capital per episode, in units of $10^{-4}$, mean over 8 seeds with t-based 95\% CIs. Left: endpoints of each swept axis (6 grid points per axis); every axis is unprofitable everywhere, so no profitability boundary exists to report. Right: the full size-response grid in the push weight.}
\label{tab:f5}
\small
\begin{tabular}{lrrl}
\toprule
Axis & Low end & High end & Boundary \\
\midrule
Kyle $\lambda$ (0 to 0.8) & $-1.3$ & $-31.7\ [-44.3, -19.2]$ & none \\
$\eta$ (0 to 0.8) & $-2.2$ & $-22.4\ [-23.5, -21.4]$ & none \\
Follower gain (0 to 120) & $-2.6$ & $-6.7\ [-10.1, -3.3]$ & none \\
\bottomrule
\end{tabular}
\hspace{1.2em}
\begin{tabular}{lrl}
\toprule
Push wt. & PnL & 95\% CI \\
\midrule
0.10 & $-0.14$ & $[-0.25, -0.03]$ \\
0.25 & $-0.50$ & $[-0.77, -0.24]$ \\
0.50 & $-1.53$ & $[-2.05, -1.00]$ \\
0.75 & $-3.07$ & $[-3.84, -2.30]$ \\
1.00 & $-5.14$ & $[-6.16, -4.12]$ \\
\bottomrule
\end{tabular}
\end{table}

The result (\cref{tab:f5}, \cref{fig:f5}) is the outcome the theory above predicts. On every axis the pump loses money at every sampled coefficient, with every per-point 95\% CI excluding zero, including the realistic center of the grid ($\lambda = 0.1$, $\eta = 0.05$: impact PnL $-3.4 \times 10^{-4}$, CI $[-4.3, -2.6] \times 10^{-4}$). The size response is bounded and decreasing: the least-bad size is the smallest sampled push weight and the loss grows to $-5.1 \times 10^{-4}$ at full size. A profitable region here would contradict the linear specification; the sweep instead agrees with Huberman-Stanzl at every point. The two extensions that follow settle what this sweep cannot. The concave ablation keeps the symmetric schedule and still finds no profit, which is a measurement rather than a corollary; the positive control then shows the same probe finding a profitable asymmetric round trip under concave impact and none under linear, so the linear null above is the theorem confirmed by an instrument that can fail, not one that cannot. The module's disclaimer remains explicit: it diagnoses the simulator, not a trading strategy.

\begin{figure}[t]
\centering
\begin{subfigure}[t]{0.55\linewidth}
\includegraphics[width=\linewidth]{f5-boundaries}
\caption{Boundary sweeps per impact axis.}
\end{subfigure}\hfill
\begin{subfigure}[t]{0.42\linewidth}
\includegraphics[width=\linewidth]{f5-size-response}
\caption{Size response in the push weight.}
\end{subfigure}
\caption{The manipulation probe: impact-attributed PnL is negative everywhere and worsens with impact strength and with size.}
\label{fig:f5}
\end{figure}

\subsection{Normalized-flow concavity: a falsifiable symmetric probe}\label{sec:f5-concave}

The linear F5 sweep is primarily a consistency check under the model assumptions of \citet{huberman2004price,gatheral2010no}. To make the probe capable of returning a positive, the engine supplies an opt-in permanent-impact exponent on dimensionless crowd flow $x=Q/V$:
\[
I_\kappa(Q;V)=\lambda\,\operatorname{sign}(x)|x|^\kappa,
\qquad x=Q/V.
\]
The temporary own-size term remains $\eta q_i/V$. This normalization is load-bearing: applying a power to raw shares and dividing by $V$ afterwards would change units with $\kappa$ and make results depend on the arbitrary share denomination. A regression test rescales capital, flow and $V$ together and requires the nonlinear impact fraction to remain unchanged. At $\kappa=1$ the engine uses the original expression $(\lambda Q+\eta q_i)/V$ exactly, preserving every linear golden bit; only $\kappa\ne1$ enters the \texttt{powf} path outside the cross-runtime golden guarantee.

We rerun the symmetric push-and-unwind sweep at $\kappa\in\{0.5,0.7\}$, eight seeds per point, with paired zero-impact references and pointwise $t$ intervals. No sampled mean is positive. At the canonical calibration, impact-attributed P\&L is $-8.9\times10^{-4}$ (95\% CI $[-19.6,1.8]\times10^{-4}$) for $\kappa=0.5$, $-13.5\times10^{-4}$ (CI $[-21.2,-5.9]\times10^{-4}$) for $\kappa=0.7$, and $-3.4\times10^{-4}$ (CI $[-4.3,-2.6]\times10^{-4}$) for the linear arm. For $\kappa=0.5$, 7 of 23 stored sweep slots have intervals crossing zero; repeated canonical slots reduce these to five distinct configurations. At $\kappa=0.7$, 0 of 23 cross zero. These are finite-grid, pointwise summaries of this symmetric schedule, not a universal no-manipulation result.

\begin{figure}[t]
\centering
\includegraphics[width=0.72\linewidth]{f5-concave}
\caption{Pointwise impact-attributed P\&L intervals for the sampled symmetric schedule under linear and normalized-flow concave impact.}
\label{fig:f5-concave}
\end{figure}

The null identifies the schedule's limitation rather than closing the economic question: profitable constructions under nonlinear permanent impact can be asymmetric. The exploratory family in \cref{sec:f5-positive-control} searches such shapes and reports selection-adjusted inference.


\subsection{Exploratory asymmetric-schedule positive control}\label{sec:f5-positive-control}

Concavity makes slicing matter. For positive dimensionless flow $x=Q/V$, the sum of the $n$ per-slice impact increments before the engine's multiplicative price compounding is
\[
n I_\kappa(Q/n;V)=n^{1-\kappa} I_\kappa(Q;V),
\]
so a slowly accumulated leg and a block liquidation need not cancel when $\kappa<1$. The identity is about the impact function, not the compounded terminal price. It motivates a positive control; it does not guarantee that one of our finite schedules must profit after compounding, temporary costs and follower flow.

The search crosses five up/down duration ratios, three block fractions, three exponents and three market arms: 135 cells, each on eight seeds. All displayed cell intervals are pointwise. Because the reported winner is selected after seeing the grid, its headline inference uses a two-sided Bonferroni familywise 95\% interval over all 135 cells (Student $t$, 7 degrees of freedom); the evidence stores the family size, critical value and adjusted interval.

The best observed cell is the no-follower, zero-temporary-impact arm at $\kappa=0.5$: nine bars of uniform accumulation followed by one-bar liquidation yield $+21.8\times10^{-4}$. Its pointwise 95\% interval is $[17.4,26.2]\times10^{-4}$ and its familywise interval remains positive at $[10.0,33.6]\times10^{-4}$. The same sampled schedule with temporary impact yields $+18.4\times10^{-4}$, with familywise interval $[6.4,30.4]\times10^{-4}$. The canonical follower arm has no pointwise-positive cell. At $\kappa=0.7$ no sampled mean is positive. Among the 45 sampled linear cells, every mean is negative; the $1{:}9$ and $9{:}1$ theory-arm values are numerically near-symmetric, not bit-identical. These statements concern the sampled grid only.

\begin{figure}[t]
\centering
\includegraphics[width=0.9\linewidth]{f5-positive-control}
\caption{Exploratory asymmetric-schedule search. Shading shows pointwise 95\% intervals; selection-adjusted intervals for the two selected $\kappa=0.5$ winners are reported in the text. Two sampled arms, not one, retain positive familywise intervals.}
\label{fig:f5-positive-control}
\end{figure}

The positive control shows that this instrument can detect a profitable sampled round trip under normalized concave impact. It does not establish a universal region, an optimum, or a theorem-level correspondence: the search fixes push size, $\lambda$, $V$, total leg length and most follower gains, and omits short-side, overshooting and longer-hold trips. Conversely, the absence of a positive among 45 sampled linear cells is consistent with the cited linear-impact model under its assumptions, not empirical proof about every configuration of the product.


\subsection{F6: Adverse-selection markout decomposition}\label{sec:f6}

\texttt{compare\_informed\_vs\_uninformed} runs twenty-four paired episodes in which the identical meta-order flow is sided on the alpha in one leg and on a coin flip in the other, holding schedule, sizes and price path fixed, and reports mean maker markout per filled unit at horizons of 1, 5 and 20 bars. Markout is measured in the tape's absolute price units (the mid starts at 100) per unit of filled quantity; the evidence commits per-episode vectors, and the gap CIs below are t-based 95\% over the twenty-four paired episodes.

\begin{table}[t]
\centering
\caption{Maker markout per filled unit (price units per unit quantity) against informed versus paired-uninformed flow, 24 paired episodes; gap CIs are t-based 95\% over per-episode paired gaps. The gap is the price of trading against information and grows with horizon; the levels stay positive at every horizon, a separate finding discussed in the text.}
\label{tab:f6}
\small
\begin{tabular}{lrrrl}
\toprule
Horizon (bars) & Informed & Uninformed & Gap & Gap 95\% CI \\
\midrule
1 & 0.689 & 0.761 & 0.072 & $[0.052, 0.094]$ \\
5 & 0.489 & 0.791 & 0.302 & $[0.228, 0.383]$ \\
20 & 0.386 & 0.823 & 0.436 & $[0.303, 0.579]$ \\
\bottomrule
\end{tabular}
\end{table}

Two certifications must be separated here: the sign of the \emph{gap} and the sign of the \emph{levels}, and they come out on opposite sides.

The gap certification passes at every horizon (\cref{tab:f6}, \cref{fig:f6}): makers filled by informed flow mark out worse than the paired-uninformed control at $h = 1$, 5 and 20, the per-unit gap widens from 0.072 to 0.436 as the information realizes, and every horizon's paired CI excludes zero, including the smallest ($h = 1$: $[0.052, 0.094]$ over per-episode paired gaps). That is the signature shape of adverse selection: the loss is in the drift after the fill, not in the fill itself. The claim is scoped to the gap, the differential cost of informed over uninformed flow, and not to maker profitability levels.

The level result goes against the environment, and we report it with the same candor as F4. Makers earn a \emph{positive} mean markout against informed flow at every horizon (0.689, 0.489, 0.386 per unit), which is on the wrong side of the \cref{sec:principles} standard that a realistic market should not let makers profit from informed flow. The cause is calibration and design, not accounting: fills are struck at the pre-move mid, so every fill books the full quoted depth as spread capture, the default quote depths are wide relative to the drift the $\alpha = 2.0$ meta-orders realize over 20 bars, and the price path is exogenous, so quotes never widen in equilibrium response to toxicity as they would in \citet{glosten1985bid}. At this calibration the probe therefore measures markout accounting against a fixed drift, in which informed flow costs makers relative to uninformed flow but not in absolute terms. Recalibrating the informed alpha or striking fills at post-impact prices so the level turns negative, and adding a variant where quotes widen with realized toxicity, are the named follow-ups.

The single-episode maker decomposition illustrates the exact identity on one seeded episode (seed 0); it is an illustration, not an aggregate finding. The wider quoter (\texttt{maker\_0}, 135 fills) keeps a positive markout at $h = 20$ (0.526 per unit, toxic-fill rate 0.296): the flow is not selecting against it. The tighter quoter (\texttt{maker\_1}, 92 fills) shows the subsidized pattern: aggregate markout is still positive at $h = 20$ (0.197 per unit) but its meta-order fills alone lose 0.197 per unit on 62\% of its filled quantity, with a toxic-fill rate of 0.533 at $h = 20$, so noise flow is paying for the informed flow. Spread capture and adverse drift decompose exactly ($\texttt{spread\_capture} + \texttt{adverse\_drift} = \texttt{markout}$, in price units times quantity): \texttt{maker\_1} captures 82.8 of spread and gives back 64.7 to adverse drift by $h = 20$.

\begin{figure}[t]
\centering
\includegraphics[width=0.7\linewidth]{f6-markouts}
\caption{Markout per filled unit by horizon: informed flow versus the paired-uninformed control.}
\label{fig:f6}
\end{figure}

\subsection{F7: Failure-mode distributions}\label{sec:f7}

Every rollout of the eight-policy field (the six reference policies plus the two behavioral counterparties), across all three tiers and sixteen seeds (384 episodes), is classified by \texttt{classify\_episode\_failure} into the worst-case-first taxonomy: cascade-wiped, bankrupt, stopped-out, mandate breach (structural, drawdown or inventory), or clean, with a per-scenario sampled mandate and a 50\% drawdown stop in force.

\begin{table}[t]
\centering
\caption{Failure-mode counts per tier, 128 episodes each (8 policies $\times$ 16 seeds). No episode reaches bankruptcy or a cascade wipe; failure is dominated by mandate breaches, and drawdown breaches triple from Calm to Extreme.}
\label{tab:f7}
\small
\setlength{\tabcolsep}{5pt}
\begin{tabular}{lrrrrrrr}
\toprule
Tier & Clean & Stopped out & Mandate struct. & Mandate DD & Mandate inv. & Bankrupt & Clean rate \\
\midrule
Calm & 72 & 0 & 36 & 11 & 9 & 0 & 0.56 \\
Hard & 68 & 0 & 36 & 14 & 10 & 0 & 0.53 \\
Extreme & 50 & 5 & 31 & 32 & 10 & 0 & 0.39 \\
\bottomrule
\end{tabular}
\end{table}

One mechanism drives the largest counts and must be stated first: mandates are sampled per scenario and assigned to fixed policies that cannot observe or adapt to them, so a structural breach is guaranteed whenever the draw pairs, say, a shorting policy with a long-only mandate. The flat structural-breach counts across tiers (36, 36, 31) are the signature of that assignment mechanism, not a behavioral finding about agents; a mandate-aware field could invert the distribution entirely, and measuring one is the interesting follow-up.

\Cref{tab:f7} and \cref{fig:f7} show the shape of failure as stress rises. The clean rate falls from 0.56 on Calm to 0.39 on Extreme. Structural mandate breaches are roughly flat across tiers, as the mechanism above dictates: a shorting policy under a sampled long-only mandate breaks the rule regardless of the weather. The stress-sensitive modes are the drawdown family: mandate drawdown breaches triple from 11 on Calm to 32 on Extreme, and hard stop-outs at the 50\% line appear only on Extreme (5 episodes, four of them \texttt{max\_sharpe}). No policy in this field ever reaches bankruptcy or a cascade wipe; the taxonomy's worst states are reachable but not by these baselines at this stop level. Per policy, \texttt{flat} is clean 48/48 by construction, while \texttt{overconfident} is the worst citizen on Calm (7 of its 16 episodes end in a drawdown breach) and \texttt{max\_sharpe} is the worst on Extreme (3 clean of 16). The headline, scoped to what was measured: on this mandate-blind reference field, half of everything that goes wrong is process, not PnL, which is exactly the half a return-ranked board never sees. Whether the same holds for agents that can read their mandate is an open measurement, not a claim of this paper.

\begin{figure}[t]
\centering
\includegraphics[width=\linewidth]{f7-failures}
\caption{Failure-mode distribution per tier per policy over 384 classified episodes.}
\label{fig:f7}
\end{figure}

\subsection{F8: Ecology outcomes under shocks, replicated over seeds}\label{sec:f8}

\texttt{run\_ecology} seeds the replicator with the eight baseline species over a shared-book market via \texttt{market\_payoffs} (12 generations, field size 8, innovation every 4 generations breeding parameter-perturbed variants of the leader) and runs trajectories under a steady all-Calm control schedule and a \texttt{regime\_shocks} schedule rotating Calm (generations 0 to 3), Hard (4 to 7) and Extreme (8 to 11). The probe is replicated over eight replicator seeds per schedule; \cref{tab:f8} reports the cross-seed winner distribution, and the seed-0 trajectories (\cref{fig:f8}) remain the illustrative detail. Note the two schedules chain episode seeds differently (the shock schedule strides seeds across generations), so control and shocked runs at the same replicator seed are matched in configuration, not in price paths.

The replicator flow rule sits in the market-ecology tradition: strategies as species whose profitability decays as their share of the flow grows \citep{farmer2002market}, and populations whose composition, not any single agent's optimization, drives aggregate dynamics \citep{lux1999scaling}. The implementation abstracts away much of what those models study: capital flows react instantly (no fund-flow lags), there are no leverage or margin constraints beyond the strategies' own sizing, and there is no entry beyond the single leader-perturbation innovator. Its outputs are trajectories of this replicator map, not claims about market ecology in general.

\begin{table}[t]
\centering
\caption{Replicated ecology outcomes over 8 replicator seeds per schedule: distribution of the final dominant lineage (bred variants credited to their founding species). The root winner differs between the schedules on 5 of 8 seeds; the seed-0 replacement of the volatility targeter by the unlevered long book occurs on 1 of 8 seeds.}
\label{tab:f8}
\footnotesize
\setlength{\tabcolsep}{4pt}
\begin{tabular}{lrr}
\toprule
Final dominant lineage & Control (all Calm) & Shocked (Calm/Hard/Extreme) \\
\midrule
kelly\_vol\_target & 4 & 7 \\
equal\_weight\_long & 2 & 1 \\
flat (no resolved dominant) & 2 & 0 \\
\bottomrule
\end{tabular}
\end{table}

The replication overturns the single-seed narrative, and we report that plainly. At seed 0 the story was clean: \texttt{kelly\_vol\_target} sweeps the control board (peak share 1.00, final 0.81), collapses under the shock schedule once the regime turns Hard, and \texttt{equal\_weight\_long} inherits dominance (final 0.67). Across eight seeds that story is the exception. Under the shock schedule a \texttt{kelly\_vol\_target} lineage ends dominant on 7 of 8 seeds (final shares near 0.90); the seed-0 handoff to the long book happens once. The control side is itself seed-dependent: \texttt{kelly\_vol\_target} wins 4 seeds, \texttt{equal\_weight\_long} wins 2, and on 2 seeds no species resolves dominance at all (all eleven end persistent at small shares, with a \texttt{flat} variant holding the largest share at 0.10). The root-level winner differs between the two schedules on 5 of 8 seeds, so shocks do reorder outcomes more often than not, but with no consistent direction. Extinction pressure is more stable than identity: shocked runs produce a resolved dominant on 8/8 seeds with 4 to 7 extinctions per run (at seed 0, both behavioral counterparties are gone or collapsed on both schedules, \texttt{overconfident} by generation 1).

The honest conclusion is methodological. One replicator trajectory is one draw from a heavy-tailed outcome distribution, and single-run ecology narratives, including the one this paper's earlier draft promoted to the abstract, do not survive replication. What the probe reliably produces at this scale is the outcome distribution itself: which lineages can win, how often shocks reorder the board, and how much of the population goes extinct, per committed seed set. Sharper claims need more seeds, varied shock orderings and path-matched schedule pairs, all of which are mechanical extensions of the committed script.

\begin{figure}[t]
\centering
\begin{subfigure}[t]{0.49\linewidth}
\includegraphics[width=\linewidth]{f8-ecology-control}
\caption{Control: steady Calm.}
\end{subfigure}\hfill
\begin{subfigure}[t]{0.49\linewidth}
\includegraphics[width=\linewidth]{f8-ecology-shocked}
\caption{Shocked: Calm, then Hard, then Extreme.}
\end{subfigure}
\caption{Seed-0 population-share trajectories over 12 replicator generations, the illustrative single run behind \cref{tab:f8}'s distribution. At this seed the regime rotation replaces the calm-regime winner with the unlevered long book; across the eight-seed replication that replacement occurs once.}
\label{fig:f8}
\end{figure}

% Predictability probe: the in-band generator-inversion attack, measured.
% Standalone fragment; evidence in paper/evidence/predictability.json, figure in
% paper/figures/predictability.pdf, script paper/src/make-predictability.py.

\subsection{Predictability probe: measuring the in-band inversion attack}\label{sec:predictability}

\Cref{sec:limitations} names a threat the leak-free interface property does not cover: a scenario is a pure function of a public 64-bit seed under a public deterministic transform, so future bars are in principle computable from past bars, and no forbidden operation is needed. This probe converts that named threat into a measurement. Three adversaries predict the next bar's return from a causal prefix only, on the three volatility tiers, 16 seeds per tier (4 symbols, 120 days, 30 warmup bars): an \emph{unconditional baseline} that predicts the expanding prefix-mean return per symbol; an \emph{honest statistical adversary}, a ridge AR(5) refit on the prefix at every bar, with no knowledge of the generator; and a \emph{known-seed oracle}, which replays the generator from the true seed and therefore predicts every future bar exactly, the upper bound on any inversion attack. Each predictor is also run as a frictionless unit-gross long/short sign-following policy and scored through the public \texttt{score\_run} kernel; we report the kernel's deflated-Sharpe probability.

\paragraph{The honest adversary extracts almost nothing.} On Calm the ridge AR reaches $55.4\% \pm 2.8\%$ directional accuracy against the baseline's $54.1\% \pm 2.4\%$: it detects the tier's designed momentum autocorrelation, worth about one accuracy point over drift-following, and its policy DSR ($0.97 \pm 0.06$ vs.\ $0.90 \pm 0.12$) reflects the same mild, real, and disclosed structure rather than any transform inversion. On the stressed tiers the edge vanishes: $51.2\% \pm 3.3\%$ on Hard and $50.6\% \pm 2.1\%$ on Extreme, statistically indistinguishable from coin flipping, with MSE no better than the unconditional baseline (Hard: $1.24\times10^{-4}$ vs.\ $1.22\times10^{-4}$; Extreme: $1.27\times10^{-3}$ vs.\ $1.17\times10^{-3}$) and policy DSRs of $0.48 \pm 0.35$ and $0.35 \pm 0.24$. The amplified tiers behave as near-martingales to a regression adversary: in-band statistical prediction recovers the disclosed stylized structure and nothing beyond it.

\paragraph{The oracle bound is total, and a bounded seed search reaches it.} The known-seed oracle scores $100\%$ accuracy, zero MSE, and DSR $1.00$ on every tier: determinism means the gap between an honest adversary and a successful inverter is the entire distance from coin flipping to omniscience, so the question is only the cost of inversion. We implemented the strongest cheap version: a brute scan over a bounded seed band of width $2^{16}$ centered on the published \texttt{start\_level}, matching each candidate's first-bar closes against the single observed opening bar. The scan recovered the true seed in 16 of 16 trials with zero collisions, from a prefix of one bar, at $14.7\,\mu$s per candidate ($\approx 0.96$\,s per scan); every recovered seed was then verified to reproduce the full 30-bar observed prefix on its tier. The lesson is sharp: when the deployed seed band is bounded and public, as a \texttt{(start\_level, num\_levels)} spec makes it, in-band inversion is not cryptanalysis but a table scan, and any such deployment should treat its scenarios as fully known to a motivated adversary. Extrapolating the measured per-candidate cost to the unbounded $2^{64}$ space gives $\sim 8.6$ million CPU-years, so brute-force inversion of an unpublished seed drawn from the full space is infeasible by enumeration.

\begin{figure}[t]
\centering
\includegraphics[width=0.78\linewidth]{predictability}
\caption{Next-bar directional accuracy by tier. The causal AR adversary remains near the unconditional baseline; knowing or recovering the seed gives perfect prediction.}
\label{fig:predictability}
\end{figure}

\paragraph{What remains open.} The probe brackets the attack but does not close it. Between the honest regression (no edge beyond disclosed structure) and the band scan (total recovery given a bounded public band) lies the case that matters for unbounded deployments: analytic inversion of SplitMix64 through the mul/add/div/max price transform from observed reals. SplitMix64 is not a cryptographic generator; its state transition is a fixed odd-constant addition and its output mixing is invertible, so recovering the state from a handful of \emph{exact} raw outputs is straightforward. What the adversary observes here is not raw outputs but prices, each a truncated 53-bit draw passed through data-dependent arithmetic with floors, and whether that quantization and composition resists lattice or guess-and-determine state recovery is a cryptanalytic question this paper does not answer. The honest reading of the probe: the environment's unpredictability rests on the secrecy and entropy of the seed, not on any hardness of the transform, and published bounded seed bands forfeit both.


% Sealed evaluation seeds: closing the in-band table-scan hole the predictability
% probe measured. Standalone fragment; evidence in paper/evidence/sealed-seeds.json,
% script paper/src/make-sealed-seeds.py. Mechanism: sharpearena::sealed_seed
% (crates/sharpearena/src/scenario_gen.rs), Python sealed_eval_seeds /
% evaluate_eval_set(salt=...) (python/sharpearena/eval_seeds.py).

\subsection{Sealed evaluation seeds: closing the measured hole}\label{sec:sealed-seeds}

The predictability probe (\Cref{sec:predictability}) left one attack fully realized: when the deployed seed band is bounded and public, in-band inversion is a table scan, not cryptanalysis. The shipped held-out set is such a band. Its named seeds are public constants, \texttt{EVAL\_SEED\_BASE} plus fixed offsets, so anyone with the repository can enumerate them, replay every held-out scenario in advance, and present a policy that has memorized the answers. The disjointness proof (every eval seed lies at or above \texttt{EVAL\_SEED\_BASE}, every train seed below it) says nothing about this; it guarantees the agent never \emph{trained} on an eval seed, not that the eval seeds are unknown. This section adds the opt-in mechanism that closes exactly that hole, and measures the closure with the probe's own adversary.

\paragraph{Mechanism.} A sealed evaluation keeps the public band \emph{structure} and hides the concrete seeds behind a salt. The evaluator supplies a secret salt $k$ (at least 16 random bytes; the derivation is only as unguessable as the salt), and the seed for named slot $i$ is
\[
\mathrm{seed}_i \;=\; \texttt{EVAL\_SEED\_BASE} \;+\; \bigl(F(k, i) \bmod (2^{64} - \texttt{EVAL\_SEED\_BASE})\bigr),
\]
where $F$ absorbs the salt bytes and their length with FNV-1a/64 and drives the digest, a fixed domain tag, and the slot index through three rounds of the SplitMix64 finalizer, chaining the digest back in between rounds. It is built entirely from primitives the crate already carries (no cryptographic dependency was added), and it is exposed identically from Rust (\texttt{sealed\_seed(salt, slot)}), the pyo3 binding, and Python (\texttt{sealed\_eval\_seeds(salt)} returns the same \texttt{name}$\to$\texttt{seed} map as the public set; \texttt{evaluate\_eval\_set(\ldots, salt=)} runs the unchanged protocol on it). Three properties are pinned by tests on both sides of the binding: every sealed seed lands in $[\texttt{EVAL\_SEED\_BASE}, 2^{64})$, so disjointness from the train band holds by construction and is checkable by anyone \emph{without} the salt; the map is deterministic in $(k, i)$, with a cross-language pinned value; and distinct salts (down to a single flipped bit) give distinct seeds. The public set, its version string, and the committed regression snapshot are untouched; sealing is opt-in.

We are explicit about what the construction is. It is a keyed PRF-\emph{style} derivation, not a certified MAC: FNV-1a is not collision resistant and the SplitMix64 finalizer is a public bijection, so an adversary holding several revealed $(i, \mathrm{seed}_i)$ pairs may be able to recover the absorbed salt digest algebraically. What it buys is the property the probe showed was missing: with a high-entropy salt, an adversary who knows the band structure, the slot names, the derivation, and the generator cannot enumerate the seeds, because the $2^{64}$-wide band is no longer collapsed onto a $2^{16}$-wide public window. What it does not buy is any hardness of the generator itself: once a sealed scenario's tape is observed, the adversary faces the same open question as before, analytic inversion of SplitMix64 through the price transform, and a revealed salt spends every seed it derives.

\paragraph{Commit-reveal workflow.} The salt lives outside the repository and outside the evaluated agent's reach. Before the run the evaluator publishes a commitment, $\mathrm{SHA\text{-}256}(k)$, which pairs naturally with SharpeBench's forward attestation: the attestation binds the commitment before any score exists. After the run the salt is revealed; anyone can recompute the seeds from $(k, i)$, check the commitment, and replay the evaluation byte for byte through the same deterministic generator and kernel. The evaluation is therefore secret while it matters and reproducible afterwards, and a salt is single use.

\paragraph{The probe's adversary, re-run.} We re-ran the band-scan adversary of \Cref{sec:predictability} unchanged against two derivations of 16 named slots: \emph{public}, the eight committed offsets plus eight further offsets drawn uniformly from the same $2^{16}$ band (a hypothetical extension of the public set, derived the same way); and \emph{sealed}, the same 16 slots under a 32-byte salt from the operating system's entropy source, committed by hash before the run. The adversary builds the first-bar-close table for the $2^{16}$ candidates at the public band start (0.75\,s, about $11\,\mu$s per candidate), matches the single observed opening bar against it (a few milliseconds per trial), and verifies any hit against the 30-bar prefix on the Hard tier. On the public band it recovered 16 of 16 seeds with zero collisions, all 16 prefix-verified: the shipped held-out set is fully known to this adversary. On the sealed band it recovered 0 of 16; the pair the mechanism is judged by is 16/16 public against 0/16 sealed. That is not luck at the margin: the scan covers a $2^{16}/(2^{64} - 10^{6}) \approx 3.6\times10^{-15}$ fraction of the band, so the expected number of recoveries at this budget is $5.7\times10^{-14}$, and the probe's own extrapolation ($\sim 8.6$ million CPU-years for the full space) is the cost of extending the scan until it succeeds. Finally, revealing the salt replayed all 16 sealed scenarios exactly (16 of 16 first-bar matches against the recomputed seeds), which is the commit-reveal guarantee in operation. Evidence, including the commitment, the revealed salt, and every per-trial record, is in \texttt{paper/evidence/sealed-seeds.json}.

\paragraph{Scope.} This closes the specific hole the probe measured, the public bounded band, and nothing more. The honest statistical adversary of \Cref{sec:predictability} is unaffected in either direction (it never used the seed), and the open cryptanalytic question about SplitMix64 through the price transform remains open. The environment's unpredictability still rests on the secrecy and entropy of the seed; sealed seeds are the mechanism that lets an evaluator actually supply both, while keeping the disjointness proof public and the evaluation replayable.

